\documentclass{article}
\usepackage[margin=2cm]{geometry}
\usepackage{amsmath}
\usepackage{graphicx}
\usepackage{booktabs}
\usepackage[utf8]{inputenc}
\usepackage{ragged2e}
\setlength{\emergencystretch}{3em}
\begin{document}
\sloppy

result given by


\begin{equation}
\begin{array} { r l } { Z = \frac { 1 } { | \mathcal { W } | } \int _ { \mathrm { C a r t o n } } [ d \sigma ] ~ e ^ { - \frac { \beta \mathbf { g } _ { \mathrm { F M } } ^ { 2 } I } { \mathrm { T r } } ( \sigma ^ { 2 } ) } \mathrm { d e t } _ { \mathrm { A d } } \left ( \sin ( \mathrm { i } \pi \sigma ) \mathrm { e } ^ { \frac { 1 } { 4 } \mathrm { i } ( \sigma ) } \right ) } & { } \\ { \times \prod _ { i } \mathrm { d e t } _ { \mathrm { R}}}\end{array}
\end{equation}


3 the radius of S, o is a dimensionless matrix, and f is de


\begin{equation}
f ( y ) = { \frac { i \pi y ^ { 3 } } { 3 } } + y ^ { 2 } \log \left ( 1 - e ^ { - 2 \pi i y } \right ) + { \frac { i y } { \pi } } \mathrm { L i } _ { 2 } \left ( e ^ { - 2 \pi i y } \right ) + { \frac { 1 } { 2 \pi ^ { 2 } } } \mathrm { L i } _ { 3 } \left ( e ^ { - 2 \pi i y } \right ) - { \frac { \zeta ( 3 ) } { 2 \pi ^ { 2 } } } \qquad \quad ( 2 )
\end{equation}


The quotient by the Weyl group in amounts to division by a simple numerical factor |W| =
2N!. The integral over o is not restricted to a Weyl chamber. Though this localization.
result was obtained in the IR theory, it is expected to hold in the UV due to the assumed.
Q-exactness of the irrelevant UV completion terms. One may rewrite the partition function.
in terms of the free energy as


\begin{equation}
\begin{array} { c } { Z = \frac { 1 } { | \mathcal { W } | } \int _ { \mathrm { C a r t a n } } [ \mathrm { d } \sigma ] ~ e ^ { - F ( \sigma ) } + O \left ( e ^ { \frac { - 4 \theta s ^ { 3 } x } { g _ { Y M } ^ { 2 } } } \right ) } \\ { F ( \sigma ) = \frac { 4 \pi ^ { 3 } r } { g _ { Y M } ^ { 2 } } \mathrm { T r } \sigma ^ { 2 } + \mathrm { T r } _ { \Lambda } F _ { V } ( \sigma ) + \sum _ { I } \mathrm {}}\end{array}
\end{equation}


The definitions of Fv() and FH() follow simply from , and using one may obtain the.
following large argument expansions


\begin{equation}
F _ { V } ( \sigma ) \approx \frac { \pi } { 6 } | \sigma | ^ { 3 } - \pi | \sigma | \qquad \qquad F _ { H } ( \sigma ) \approx - \frac { \pi } { 6 } | \sigma | ^ { 3 } - \frac { \pi } { 8 } | \sigma | \qquad \qquad ( 4 )
\end{equation}


It was argued in that in the large N limit, the perturbative Yang-Mills termi.e. the first
term in the expression for F(o) incan be neglected, as can be the instanton contributions.
Thus in our evaluation of the free energy, we will only concern ourselves with the contributions coming from Fv(o) and Fh(). The first step in the evaluation of is recasting the.
matrix integral in a simpler form. The integral over o in is an integration over the Coulomb
branch, which is parameterized by the non-zero vevs of o. One may write


\begin{equation}
{ \boldsymbol \sigma } = \mathrm { d i a g } \{ \lambda _ { 1 }, \dots, \lambda _ { N }, - \lambda _ { 1 }, \dots, - \lambda _ { N } \} \qquad ( 5 )
\end{equation}


since USp(2N) has N elements in its Cartan. The integration variables are these N A..
Normalizing the weights of the fundamental representation of USp(2N) to be e; with ej
forming a basis of unit vectors for R, it follows that the adjoint representation has weights.
2e, and e,  e, for all i  j, whereas the anti-symmetric representation has only weights

\end{document}
